% GENERATED FILE. Edit canonical lessons, then run npm run generate:books.
% canonical-source-hash: fnv1a64:90b438560aaeb8d6
% canonical-lessons: FR-C340-fiable, FR-C340-credible, FR-C340-douteux, FR-C340-vraisemblable, FR-C340-invraisemblable

\chapter{Croire ou pas}
\label{ch:fr-croire-ou-pas}
\hlchaptermodality{340}

\begin{chapteropening}
\textbf{By the end of this chapter:} \emph{I can say whether a person or a story is reliable, believable, dubious, likely or hard to believe.}

Complete the last lesson of chapter 340: I can say whether a person or a story is reliable, believable, dubious, likely or hard to believe.
\end{chapteropening}

\section[fiable]{fiable — reliable}

\label{lesson:FR-C340-fiable}



\begin{quote}
Before the new one: say the French for indisputable, then the French for clear.
\end{quote}

\subsection*{You'll want to know: fiable}
\textbf{fiable} — \textquotedblleft{}reliable\textquotedblright{}.

\begin{quote}
Cette voiture est vieille, mais très fiable.
\end{quote}

\begin{quote}
\emph{This car is old, but very reliable.}
\end{quote}

For people and things: \emph{une source fiable}.

\begin{cousinweb}
\emph{fiable} is built on \emph{se fier} \textquotedblleft{}to trust\textquotedblright{} (EtymDB), which is usually traced to a Vulgar Latin \emph{fīdāre}, formed on Latin \emph{fīdere} \textquotedblleft{}to trust\textquotedblright{} (Lewis \& Short).
\end{cousinweb}

\begin{grammarlens}[title={What you've built}]
One more word for telling, explaining and describing at B1.
\end{grammarlens}

\subsection*{Your turn}
\begin{itemize}
\raggedright
  \item \emph{Say it:} \emph{fiable}
  \item \emph{Say it:} \emph{fiable}, once more
  \item \emph{Say it:} a sentence of your own with \emph{fiable}
  \item \emph{Recall:} say the French for indisputable, then the French for clear, then say \emph{fiable} again
\end{itemize}

\begin{tcolorbox}[breakable,colback=teal!4,colframe=teal!35!black,title={Before you move on}]
What does fiable mean? (\textquotedblleft{}Reliable\textquotedblright{}.) Say it once more.
\end{tcolorbox}
\section[crédible]{crédible — believable, credible}

\label{lesson:FR-C340-credible}



\begin{quote}
Before the new one: say the French for clear, then the French for reliable.
\end{quote}

\subsection*{You'll want to know: crédible}
\textbf{crédible} — \textquotedblleft{}believable, credible\textquotedblright{}.

\begin{quote}
Son histoire n'est pas très crédible.
\end{quote}

\begin{quote}
\emph{His story is not very believable.}
\end{quote}

One form for both genders. Close to \emph{croyable}, but more formal.

\begin{cousinweb}
\emph{crédible} is a learned borrowing from Latin \emph{crēdibilis} \textquotedblleft{}worthy of belief\textquotedblright{} (Lewis \& Short), from \emph{crēdere} \textquotedblleft{}to believe\textquotedblright{}.
\end{cousinweb}

\begin{grammarlens}[title={What you've built}]
One more word for telling, explaining and describing at B1.
\end{grammarlens}

\subsection*{Your turn}
\begin{itemize}
\raggedright
  \item \emph{Say it:} \emph{crédible}
  \item \emph{Say it:} \emph{crédible}, once more
  \item \emph{Say it:} a sentence of your own with \emph{crédible}
  \item \emph{Recall:} say the French for clear, then the French for reliable, then say \emph{crédible} again
\end{itemize}

\begin{tcolorbox}[breakable,colback=teal!4,colframe=teal!35!black,title={Before you move on}]
What does crédible mean? (\textquotedblleft{}Believable, credible\textquotedblright{}.) Say it once more.
\end{tcolorbox}
\section[douteux]{douteux — doubtful, dubious}

\label{lesson:FR-C340-douteux}



\begin{quote}
Before the new one: say the French for reliable, then the French for believable.
\end{quote}

\subsection*{You'll want to know: douteux}
\textbf{douteux} — \textquotedblleft{}doubtful, dubious\textquotedblright{}.

\begin{quote}
Ce restaurant a l'air douteux : il n'y a personne.
\end{quote}

\begin{quote}
\emph{This restaurant looks dubious: there is nobody there.}
\end{quote}

Feminine \emph{douteuse}. \emph{Il est douteux que} takes the subjunctive.

\begin{cousinweb}
\emph{douteux} is built from \emph{doute} \textquotedblleft{}doubt\textquotedblright{} and the suffix \emph{-eux} (EtymDB). \emph{doute} comes from \emph{douter}, inherited from Latin \emph{dubitāre} \textquotedblleft{}to waver, to doubt\textquotedblright{} (Lewis \& Short).
\end{cousinweb}

\begin{grammarlens}[title={What you've built}]
One more word for telling, explaining and describing at B1.
\end{grammarlens}

\subsection*{Your turn}
\begin{itemize}
\raggedright
  \item \emph{Say it:} \emph{douteux}
  \item \emph{Say it:} \emph{douteux}, once more
  \item \emph{Say it:} a sentence of your own with \emph{douteux}
  \item \emph{Recall:} say the French for reliable, then the French for believable, then say \emph{douteux} again
\end{itemize}

\begin{tcolorbox}[breakable,colback=teal!4,colframe=teal!35!black,title={Before you move on}]
What does douteux mean? (\textquotedblleft{}Doubtful, dubious\textquotedblright{}.) Say it once more.
\end{tcolorbox}
\section[vraisemblable]{vraisemblable — likely, believable}

\label{lesson:FR-C340-vraisemblable}



\begin{quote}
Before the new one: say the French for believable, then the French for doubtful.
\end{quote}

\subsection*{You'll want to know: vraisemblable}
\textbf{vraisemblable} — \textquotedblleft{}likely, believable\textquotedblright{}.

\begin{quote}
C'est une explication vraisemblable.
\end{quote}

\begin{quote}
\emph{It is a likely explanation.}
\end{quote}

Literally like the truth: \emph{vrai} + \emph{semblable}.

\begin{cousinweb}
\emph{vraisemblable} is built from \emph{vrai} \textquotedblleft{}true\textquotedblright{} and \emph{semblable} \textquotedblleft{}like\textquotedblright{}, on the model of Latin \emph{vērī similis} \textquotedblleft{}like the truth\textquotedblright{} (Lewis \& Short).
\end{cousinweb}

\begin{grammarlens}[title={What you've built}]
One more word for telling, explaining and describing at B1.
\end{grammarlens}

\subsection*{Your turn}
\begin{itemize}
\raggedright
  \item \emph{Say it:} \emph{vraisemblable}
  \item \emph{Say it:} \emph{vraisemblable}, once more
  \item \emph{Say it:} a sentence of your own with \emph{vraisemblable}
  \item \emph{Recall:} say the French for believable, then the French for doubtful, then say \emph{vraisemblable} again
\end{itemize}

\begin{tcolorbox}[breakable,colback=teal!4,colframe=teal!35!black,title={Before you move on}]
What does vraisemblable mean? (\textquotedblleft{}Likely, believable\textquotedblright{}.) Say it once more.
\end{tcolorbox}
\section[invraisemblable]{invraisemblable — unlikely, hard to believe}

\label{lesson:FR-C340-invraisemblable}



\begin{quote}
Before the new one: say the French for doubtful, then the French for likely.
\end{quote}

\subsection*{You'll want to know: invraisemblable}
\textbf{invraisemblable} — \textquotedblleft{}unlikely, hard to believe\textquotedblright{}.

\begin{quote}
Il m'a raconté une histoire invraisemblable.
\end{quote}

\begin{quote}
\emph{He told me a story that was hard to believe.}
\end{quote}

With \emph{in-}, the opposite of \emph{vraisemblable}: so odd it is hard to believe.

\begin{cousinweb}
\emph{invraisemblable} is built from \emph{in-} \textquotedblleft{}not\textquotedblright{} and \emph{vraisemblable}, which is formed from \emph{vrai} \textquotedblleft{}true\textquotedblright{} and \emph{semblable} \textquotedblleft{}like\textquotedblright{}, on the model of Latin \emph{vērī similis} \textquotedblleft{}like the truth\textquotedblright{} (Lewis \& Short).
\end{cousinweb}

\begin{grammarlens}[title={What you've built}]
One more word for telling, explaining and describing at B1.
\end{grammarlens}

\subsection*{Your turn}
\begin{itemize}
\raggedright
  \item \emph{Say it:} \emph{invraisemblable}
  \item \emph{Say it:} \emph{invraisemblable}, once more
  \item \emph{Say it:} a sentence of your own with \emph{invraisemblable}
  \item \emph{Recall:} say the French for doubtful, then the French for likely, then say \emph{invraisemblable} again
\end{itemize}

\begin{tcolorbox}[breakable,colback=teal!4,colframe=teal!35!black,title={Before you move on}]
What does invraisemblable mean? (\textquotedblleft{}Unlikely, hard to believe\textquotedblright{}.) Say it once more.
\end{tcolorbox}
